\chapter*{前言}
\addcontentsline{toc}{chapter}{前言}
\setchapterheader{前言}
\markboth{前言}{前言}

有时我们手上只有一台笔记本，希望在虚拟机里准备一套 Ubuntu 环境来连接 Bumi；也有可能机器人已经有一块算力板，但仍想在隔离的虚拟机中做实验。这两种做法都会遇到同一个问题：虚拟机默认走 NAT 或共享网络，地址落在 \code{10.211.55.x} 这类网段，而 Bumi 运控板只接受 \code{192.168.55.0/24} 上的直连通信。

因此，本手册只讲一件事：先把 Ubuntu 装进虚拟机，再把虚拟机的网络从共享模式改成桥接模式，让虚拟机像一台独立的物理机一样接入机器人所在的网段，并最终完成 DDS 通信验证。操作环境分两种：Windows x86\_64 主机使用 VMware 与 amd64 镜像，MacOS（Apple Silicon）主机使用 Parallels Desktop 及其提供的 arm 镜像。

内容按“先看懂再动手”的顺序组织。第零章安装虚拟机并准备好 Ubuntu 系统；第一章说明目标地址、网卡与整体链路；第二章把虚拟机网络改成桥接；第三章配置静态地址；第四章放行 DDS 端口并测试连通；第五章运行程序验证通信。附录提供常见故障速查和恢复原配置的方法。

每一步都给出操作命令、预期结果和常见问题。命令中的网卡名 \code{enp0s5}、地址 \code{192.168.55.101} 请按自己的实际环境替换：网卡名用 \code{ip -brief addr show} 查看，地址不能与运控板或其他设备冲突。

调试时请全程使用机器人吊架，确保机器人在切换控制模式前得到稳定支撑。
